\clearpage\pagestyle{otMain}\setlength{\otBlocEspacement}{2.5mm}\fontsize{11}{15}\selectfont\corrigesChapitre{16}{Translations, agrandissements et réductions}
\begin{blocLivre}{cor-16-1}\begin{corrige}{1}{Décrire le glissement}\demarche{Reproduire le déplacement horizontal et vertical.}
Le déplacement ajoute 3 à l’abscisse et 1 à l’ordonnée. L’image est donc $\res{A'(4;3)}$. Le même comptage s’applique à chaque point.
\label{fig-cor-16-1}\figureGeo[corrige]{trans-01}
\end{corrige}\end{blocLivre}
\begin{blocLivre}{cor-16-2}\begin{corrige}{2}{Construire un triangle image}\demarche{Appliquer le même déplacement aux trois sommets.}
On obtient $\res{A'(4;3),\ B'(6;3),\ C'(5;5)}$. Placer les trois images, puis relier dans le même ordre. Les côtés correspondants restent parallèles et de même longueur.
\label{fig-cor-16-2}\figureGeo[corrige]{trans-02}
\end{corrige}\end{blocLivre}
\begin{blocLivre}{cor-16-3}\begin{corrige}{3}{Un déplacement négatif}\demarche{Relier le sens du déplacement au signe du calcul.}
Vers la gauche, on soustrait 4 à l’abscisse ; vers le bas, on soustrait 1 à l’ordonnée. Donc $\res{B'(-2;2)}$.
\label{fig-cor-16-3}\figureGeo[corrige]{trans-03}
\end{corrige}\end{blocLivre}
\begin{blocLivre}{cor-16-4}\begin{corrige}{4}{Conserver des mesures}\demarche{Citer la conservation adaptée à chaque mesure.}
Une translation conserve longueurs, angles et aires. Donc $\res{A'B'=5\ \mathrm{cm}}$, $\res{\widehat{B'A'C'}=42^\circ}$ et l’aire image vaut $\res{9\ \mathrm{cm}^2}$.
\label{fig-cor-16-4}\figureGeo[corrige]{trans-04}
\end{corrige}\end{blocLivre}
\begin{blocLivre}{cor-16-5}\begin{corrige}{5}{Reconnaître une frise}\demarche{Compter les intervalles entre les motifs.}
Du premier au suivant : translation de 4 cm vers la droite. Du premier au cinquième, on compte quatre intervalles : $4\times4=\res{16}$ cm. Cinq motifs ne font pas cinq écarts.
\label{fig-cor-16-5}\figureGeo[corrige]{trans-05}
\end{corrige}\end{blocLivre}
\begin{blocLivre}{cor-16-6}\begin{corrige}{6}{Un parallélogramme}\demarche{Respecter l’ordre des sommets du quadrilatère.}
Les segments AB et A’B’ sont parallèles et de même longueur, car ils sont images par translation. Dans l’ordre A-B-B’-A’, ce sont deux côtés opposés. Un quadrilatère ayant deux côtés opposés parallèles et de même longueur est un parallélogramme.
\label{fig-cor-16-6}\figureGeo[corrige]{trans-06}
\end{corrige}\end{blocLivre}
\begin{blocLivre}{cor-16-7}\begin{corrige}{7}{Agrandir une figure}\demarche{Multiplier toutes les longueurs par le même rapport.}
Les dimensions deviennent $6\times1{,}5=\res9$ cm et $4\times1{,}5=\res6$ cm. Le périmètre vaut $2(9+6)=\res{30}$ cm, soit 1,5 fois le périmètre initial 20 cm.
\label{fig-cor-16-7}\figureGeo[corrige]{trans-07}
\end{corrige}\end{blocLivre}
\begin{blocLivre}{cor-16-8}\begin{corrige}{8}{Réduire un triangle}\demarche{Appliquer le rapport à chaque côté.}
On multiplie chaque côté par $2/5$. Les nouvelles longueurs sont $\res2$ cm, $\res{2{,}8}$ cm et $\res{3{,}2}$ cm. Les angles restent les mêmes ; les longueurs changent.
\label{fig-cor-16-8}\figureGeo[corrige]{trans-08}
\end{corrige}\end{blocLivre}
\begin{blocLivre}{cor-16-9}\begin{corrige}{9}{Une aire agrandie}\demarche{Mettre au carré le rapport des longueurs.}
Le rapport des longueurs est 3, donc celui des aires est $3^2=9$. La nouvelle aire vaut $18\times9=\res{162\ \mathrm{cm}^2}$, et non 54.
\label{fig-cor-16-9}\figureGeo[corrige]{trans-09}
\end{corrige}\end{blocLivre}
\begin{blocLivre}{cor-16-10}\begin{corrige}{10}{Un volume réduit}\demarche{Élever le rapport au cube pour un volume.}
Toutes les dimensions sont divisées par 2 ; le volume est donc multiplié par $(1/2)^3=1/8$. Il vaut $240/8=\res{30\ \mathrm{cm}^3}$.
\label{fig-cor-16-10}\figureGeo[corrige]{trans-10}
\end{corrige}\end{blocLivre}
\begin{blocLivre}{cor-16-11}\begin{corrige}{11}{Retrouver le rapport}\demarche{Prendre la racine positive du rapport d’aires.}
$k^2=45/20=2{,}25$, donc $k=\sqrt{2{,}25}=\res{1{,}5}$. Le côté devient $8\times1{,}5=\res{12}$ cm. On choisit le rapport positif.
\label{fig-cor-16-11}\figureGeo[corrige]{trans-11}
\end{corrige}\end{blocLivre}
\begin{blocLivre}{cor-16-12}\begin{corrige}{12}{Un pavé à l’échelle}\demarche{Distinguer longueur, aire et volume.}
Les dimensions deviennent $\res{4,\ 8/3,\ 2}$ cm. La face a une aire de $(12\times8)/9=\res{32/3}$ cm$^2$. Le volume devient $(12\times8\times6)/27=\res{64/3}$ cm$^3$. Les facteurs sont respectivement $1/3$, $1/9$ et $1/27$.
\label{fig-cor-16-12}\figureGeo[corrige]{trans-12}
\end{corrige}\end{blocLivre}
\begin{blocLivre}{cor-16-13}\begin{corrige}{13}{Deux déplacements}\demarche{Additionner séparément les déplacements sur chaque axe.}
Horizontalement : $5-3=2$ carreaux à droite. Verticalement : $2+4=6$ vers le haut. Le glissement total est donc $\res{2\text{ à droite et }6\text{ en haut}}$. Les déplacements s’additionnent avec leurs sens.
\label{fig-cor-16-13}\figureGeo[corrige]{trans-13}
\end{corrige}\end{blocLivre}
\begin{blocLivre}{cor-16-14}\begin{corrige}{14}{Le bon coefficient de peinture}\demarche{Calculer la nouvelle aire avant la quantité.}
La surface est multipliée par $2{,}5^2=6{,}25$, donc elle vaut $0{,}8\times6{,}25=5$ m$^2$. La peinture nécessaire est $5/4=\res{1{,}25\ \mathrm L}$ dans le modèle sans perte. Utiliser seulement 2,5 comme facteur d’aire serait faux.
\label{fig-cor-16-14}\figureGeo[corrige]{trans-14}
\end{corrige}\end{blocLivre}
\begin{blocLivre}{cor-16-15}\begin{corrige}{15}{Un pavage par translations}\demarche{Compter les copies et convertir le déplacement en longueurs.}
Il y a $6\times4=\res{24}$ motifs, chacun d’aire 4 cm$^2$. L’aire totale vaut $\res{96\ \mathrm{cm}^2}$. Le déplacement demandé est de $\res4$ cm à droite et $\res2$ cm vers le haut. Les translations conservent la taille du motif.
\label{fig-cor-16-15}\figureGeo[corrige]{trans-15}
\end{corrige}\end{blocLivre}
